% XeLaTeX 컬러 이모지 학생 활동지
% 컴파일: xelatex worksheet.tex

\documentclass[11pt, a4paper]{article}

% 한글 지원
\usepackage{fontspec}
\usepackage{polyglossia}
\setdefaultlanguage{korean}
\setotherlanguage{english}
\setmainfont{Noto Sans CJK KR}
\setsansfont{Noto Sans CJK KR}

% 이미지 (컬러 이모지)
\usepackage{graphicx}
\graphicspath{{./emoji/}}
\newcommand{\emoji}[1]{\raisebox{-0.15em}{\includegraphics[height=1em]{#1}}}
\newcommand{\emojiL}[1]{\raisebox{-0.2em}{\includegraphics[height=1.5em]{#1}}}
\newcommand{\emojiXL}[1]{\raisebox{-0.2em}{\includegraphics[height=3em]{#1}}}

% 페이지 설정
\usepackage[margin=2cm, headheight=14pt]{geometry}
\usepackage{fancyhdr}
\pagestyle{fancy}
\fancyhf{}
\rhead{\emoji{apple} 사과 가격 예측 시뮬레이터 활동지}
\cfoot{- \thepage -}

% 색상
\usepackage{xcolor}
\definecolor{primary}{HTML}{1565C0}
\definecolor{secondary}{HTML}{00796B}
\definecolor{accent}{HTML}{F57C00}
\definecolor{success}{HTML}{2E7D32}
\definecolor{purple}{HTML}{7B1FA2}
\definecolor{lightblue}{HTML}{E3F2FD}
\definecolor{lightgreen}{HTML}{E8F5E9}
\definecolor{lightyellow}{HTML}{FFF8E1}
\definecolor{lightcyan}{HTML}{E0F2F1}
\definecolor{lightpurple}{HTML}{F3E5F5}

% 표, 박스
\usepackage{tcolorbox}
\usepackage{tabularx}
\usepackage{array}
\usepackage{colortbl}

% 기타
\usepackage{enumitem}
\usepackage{setspace}
\usepackage{amssymb}  % \square 명령어
\onehalfspacing

% 섹션 스타일
\usepackage{titlesec}
\titleformat{\section}{\Large\bfseries\color{primary}}{\thesection.}{0.5em}{}
\titleformat{\subsection}{\large\bfseries\color{secondary}}{\thesubsection}{0.5em}{}

% tcolorbox 설정
\tcbset{
    boxrule=0.5pt,
    arc=3mm,
    left=5pt,
    right=5pt,
    top=5pt,
    bottom=5pt
}

\begin{document}

% ===== 표지 =====
\begin{center}
\vspace*{0.5cm}

{\fontsize{60pt}{60pt}\selectfont \emojiXL{apple}}

\vspace{0.3cm}

{\Huge\bfseries\color{primary} 사과 가격 예측 시뮬레이터}

\vspace{0.2cm}

{\LARGE\bfseries\color{secondary} 학생 활동지}

\vspace{0.8cm}

\begin{tabular}{|p{3cm}|p{6cm}|}
\hline
\cellcolor{gray!20}\textbf{학년 / 반} & \\[0.5cm]
\hline
\cellcolor{gray!20}\textbf{이\quad\quad 름} & \\[0.5cm]
\hline
\cellcolor{gray!20}\textbf{날\quad\quad 짜} & \\[0.5cm]
\hline
\end{tabular}

\vspace{0.8cm}

\begin{tcolorbox}[colback=lightyellow, colframe=accent, title={\emojiL{books} \textbf{학습 목표}}, fonttitle=\bfseries\color{accent}]
\begin{enumerate}[leftmargin=*]
    \item 품질 점수가 사과 가격에 어떤 영향을 미치는지 이해한다.
    \item AI가 데이터로부터 규칙을 학습하는 과정을 체험한다.
    \item 학습된 모델로 새로운 사과의 가격을 예측해본다.
\end{enumerate}
\end{tcolorbox}

\end{center}

\newpage

% ===== 사용 방법 =====
\begin{tcolorbox}[colback=lightblue, colframe=primary, title={\emojiL{openbook} \textbf{활동지 사용 방법}}, fonttitle=\bfseries\color{primary}]
\begin{enumerate}[leftmargin=*]
    \item[\textcolor{primary}{\textbf{①}}] 시뮬레이터를 실행하고 선생님의 안내에 따라 활동을 진행하세요.
    \item[\textcolor{primary}{\textbf{②}}] 각 활동마다 빈칸에 관찰한 내용이나 결과를 기록하세요.
    \item[\textcolor{primary}{\textbf{③}}] 궁금한 점은 손을 들고 질문하세요!
\end{enumerate}
\vspace{0.2cm}
\colorbox{blue!15}{\emoji{bulb} \textbf{팁:} \textit{정답이 없는 질문도 있어요. 자신의 생각을 자유롭게 적어보세요!}}
\end{tcolorbox}

\vspace{0.4cm}

% ===== 1차시 =====
\begin{tcolorbox}[colback=lightblue, colframe=primary]
\begin{center}
{\Large\bfseries \emoji{bluebook} 1차시: AI 학습의 원리 이해하기}
\end{center}
\end{tcolorbox}

\vspace{0.4cm}

\subsection*{\emoji{pencil} 활동 1: 품질 점수 계산하기}

사과의 당도와 무게로 품질 점수를 계산해봅시다.

\begin{tcolorbox}[colback=lightpurple, colframe=purple, title={\emoji{triangle} \textbf{품질 점수 공식}}, fonttitle=\bfseries\color{purple}]
\texttt{당도 점수 = (당도 - 10) ÷ 8 × 50}\\
\texttt{무게 점수 = (무게 - 150) ÷ 200 × 50}\\
\texttt{\textbf{품질 점수 = 당도 점수 + 무게 점수}}
\end{tcolorbox}

\vspace{0.2cm}

\textcolor{secondary}{\textbf{▶ 다음 사과의 품질 점수를 계산해보세요:}}

\begin{center}
\begin{tabular}{|c|c|c|c|c|c|}
\hline
\rowcolor{lightblue}
\textbf{번호} & \textbf{당도(Brix)} & \textbf{무게(g)} & \textbf{당도점수} & \textbf{무게점수} & \textbf{품질점수} \\
\hline
1 & & & & & \\[0.4cm]
\hline
2 & & & & & \\[0.4cm]
\hline
3 & & & & & \\[0.4cm]
\hline
\end{tabular}
\end{center}

\vspace{0.4cm}

\subsection*{\emoji{magnify} 활동 2: 데이터 생성하고 관찰하기}

시뮬레이터에서 데이터를 생성하고 관찰해봅시다.

\vspace{0.2cm}

\textcolor{secondary}{\textbf{▶ 설정값 기록:}}

\begin{center}
\begin{tabular}{|c|c|c|}
\hline
\cellcolor{gray!20} 데이터 개수: \underline{\hspace{1.5cm}}개 & 
\cellcolor{gray!20} 노이즈: \underline{\hspace{1.5cm}}\% & 
\cellcolor{gray!20} 학습률: \underline{\hspace{1.5cm}} \\
\hline
\end{tabular}
\end{center}

\vspace{0.2cm}

\textcolor{secondary}{\textbf{▶ 산점도 관찰 (그래프에서 발견한 점을 적어보세요):}}

\begin{tcolorbox}[colback=white, colframe=gray, height=2.2cm]
\end{tcolorbox}

\newpage

% ===== 활동 3 =====
\subsection*{\emoji{bolt} 활동 3: AI 학습 실행하기}

[학습 실행] 버튼을 눌러 AI가 학습하는 과정을 관찰해봅시다.

\vspace{0.2cm}

\textcolor{secondary}{\textbf{▶ 학습 결과 기록:}}

\begin{center}
\begin{tabular}{|c|c|c|c|}
\hline
\rowcolor{lightblue}
\textbf{에포크} & \textbf{손실(MSE)} & \textbf{R² 점수} & \textbf{상태} \\
\hline
 & & & $\square$양호 $\square$보통 $\square$부족 \\[0.5cm]
\hline
\end{tabular}
\end{center}

\vspace{0.2cm}

\textcolor{secondary}{\textbf{▶ 학습된 회귀식 기록:}}

\begin{tcolorbox}[colback=lightyellow, colframe=accent]
\begin{center}
{\Large\bfseries 가격 = \underline{\hspace{2cm}} × 품질 + \underline{\hspace{2cm}}}
\end{center}
\end{tcolorbox}

\vspace{0.2cm}

\textcolor{secondary}{\textbf{▶ 이 수식이 의미하는 것은? (품질이 1점 오르면 가격은?)}}

\begin{tcolorbox}[colback=white, colframe=gray, height=1.8cm]
\end{tcolorbox}

\vspace{0.2cm}

\textcolor{secondary}{\textbf{▶ [재생] 버튼으로 학습 과정을 다시 보고 느낀 점:}}

\begin{tcolorbox}[colback=white, colframe=gray, height=1.8cm]
\end{tcolorbox}

\newpage

% ===== 2차시 =====
\begin{tcolorbox}[colback=lightcyan, colframe=secondary]
\begin{center}
{\Large\bfseries \emoji{greenbook} 2차시: AI 예측 체험하기}
\end{center}
\end{tcolorbox}

\vspace{0.4cm}

\subsection*{\emoji{ruler} 활동 4: 오차 분석하기}

[오차선 표시] 체크박스를 켜고 AI의 예측이 틀린 정도를 관찰해봅시다.

\begin{tcolorbox}[colback=lightgreen, colframe=success, title={\emoji{bulb} \textbf{오차선이란?}}, fonttitle=\bfseries\color{success}]
실제 가격(데이터 점)과 AI가 예측한 가격(회귀선) 사이의 차이를 보여주는 점선입니다.\\
오차선이 길수록 AI가 그 사과의 가격을 맞추지 못했다는 뜻이에요!
\end{tcolorbox}

\vspace{0.2cm}

\textcolor{secondary}{\textbf{▶ 오차 통계 기록:}}

\begin{center}
\begin{tabular}{|c|c|c|}
\hline
\rowcolor{lightgreen}
\textbf{평균 오차} & \textbf{최대 오차} & \textbf{MSE} \\
\hline
\underline{\hspace{2cm}} 원 & \underline{\hspace{2cm}} 원 & \underline{\hspace{2cm}} \\[0.5cm]
\hline
\end{tabular}
\end{center}

\vspace{0.2cm}

\textcolor{secondary}{\textbf{▶ AI가 왜 100\% 정확하게 맞추지 못할까요?}}

\begin{tcolorbox}[colback=white, colframe=gray, height=1.8cm]
\end{tcolorbox}

\vspace{0.4cm}

\subsection*{\emoji{crystal} 활동 5: 새로운 사과 가격 예측하기}

[추론 모드]로 전환하여 새로운 사과의 가격을 예측해봅시다.

\vspace{0.2cm}

\textcolor{secondary}{\textbf{▶ 3가지 사과의 가격을 예측하고 기록하세요:}}

\begin{center}
\begin{tabular}{|c|c|c|c|c|}
\hline
\rowcolor{lightblue}
\textbf{번호} & \textbf{당도(Brix)} & \textbf{무게(g)} & \textbf{품질점수} & \textbf{예측가격} \\
\hline
1 & & & & \\[0.4cm]
\hline
2 & & & & \\[0.4cm]
\hline
3 & & & & \\[0.4cm]
\hline
\end{tabular}
\end{center}

\newpage

% ===== 활동 6 =====
\subsection*{\emoji{microscope} 활동 6: 탐구 실험}

조건을 바꿔가며 AI 학습이 어떻게 달라지는지 실험해봅시다.

\vspace{0.2cm}

\textcolor{secondary}{\textbf{▶ 학습률 실험 (학습률만 바꾸고 나머지는 같게 유지):}}

\begin{center}
\begin{tabular}{|c|c|c|c|}
\hline
\rowcolor{lightgreen}
\textbf{학습률} & \textbf{R² 점수} & \textbf{손실 그래프 모양} & \textbf{느낀 점} \\
\hline
0.001 (작게) & & & \\[0.6cm]
\hline
0.01 (보통) & & & \\[0.6cm]
\hline
0.1 (크게) & & & \\[0.6cm]
\hline
\end{tabular}
\end{center}

\vspace{0.2cm}

\textcolor{secondary}{\textbf{▶ 학습률이 너무 크면 어떤 문제가 생기나요?}}

\begin{tcolorbox}[colback=white, colframe=gray, height=1.6cm]
\end{tcolorbox}

\vspace{0.4cm}

% ===== 종합 정리 =====
\subsection*{\emoji{memo} 종합 정리}

오늘 배운 내용을 정리해봅시다.

\vspace{0.2cm}

\textbf{1. 선형 회귀란?}
\begin{tcolorbox}[colback=white, colframe=gray, height=1.4cm]
\end{tcolorbox}

\textbf{2. AI 학습에서 가장 인상 깊었던 점은?}
\begin{tcolorbox}[colback=white, colframe=gray, height=1.4cm]
\end{tcolorbox}

\textbf{3. AI가 완벽하지 않다는 것을 어떻게 알 수 있었나요?}
\begin{tcolorbox}[colback=white, colframe=gray, height=1.4cm]
\end{tcolorbox}

\vspace{0.4cm}

% ===== 자기 평가 =====
\begin{tcolorbox}[colback=lightyellow, colframe=accent, title={\emoji{star} \textbf{자기 평가}}, fonttitle=\bfseries\color{accent}]
\begin{itemize}[leftmargin=*]
    \item[$\square$] 품질 점수 계산 방법을 이해했다
    \item[$\square$] AI가 데이터로부터 학습하는 과정을 설명할 수 있다
    \item[$\square$] 학습된 모델로 새로운 데이터를 예측할 수 있다
    \item[$\square$] AI 예측에도 오차가 있다는 것을 이해했다
\end{itemize}
\end{tcolorbox}

\end{document}
